\subsection{Other Useful Lemmas}

This section contains the remaining lemmas required for proving the main theorem (Theorem \ref{thm: main theorem}).
The following lemma calculates the gap in risk when $Z$ is substituted with $\phi: \cg{X} \rightarrow \cg{Z}$.

\begin{lem}\label{lem: gap in risk when z is replaced with phi hat}
    Let $w := (w_c,w_x,w_z) \in \R^{d+2}$ and let $\ph$ be defined in  \eqref{eqn: phi star and phi hat}. Then
    \begin{equation*}
        |\rph(w) - R(w)| \leq 2|w_z|L_{(0,1)}(\uh).
    \end{equation*}
\end{lem}

\begin{proof}
For fixed $(x,y)$ and $w$, set
$\ell(x,t,y;w)=\ln(1+\exp(-y(w_c+w_x^\top x+w_zt)))$.
Its derivative in $t$ satisfies
\[
\left|\frac{\partial}{\partial t}\ell(x,t,y;w)\right|
=|yw_z|\,[1-\sigma(y(w_c+w_x^\top x+w_zt))]\leq|w_z|.
\]
We further have $\mathbb E_{X,Z}|Z-\ph(X)|
=2\mathbb P_{X,Z}(Z\neq\ph(X))
\leq2L_{(0,1)}(\uh)$.
Therefore
\[
|\rph(w)-R(w)|
\leq\mathbb E_{X,Y,Z}
  |\ell(X,\ph(X),Y;w)-\ell(X,Z,Y;w)|
\leq2|w_z|L_{(0,1)}(\uh).
\]
\end{proof}

The next lemma establishes the effect of the regularizer on $w_z$.

\begin{lem}\label{lem: bound weights in terms of lambda}
Let $\lambda_c,\lambda_x,\lambda_z>0$ and let $\whl$ and $\wsl$
be defined in \eqref{eqn: w hat lambda} and \eqref{eqn: w lambda star}.
For either $q=\whl$ or $q=\wsl$,
\[
|q_z|\leq
\min\left\{\sqrt{\frac{\ln2}{\lambda_z}},\frac1{2\lambda_z}\right\}
\leq\frac{\ln2}{\lambda_z}.
\]
\end{lem}

\begin{proof}
Suppose $q = \whl$. Using the property that $\whl$ is an optimal regularized empirical risk minimizing weights, we have 
\[
\lambda_z|\whlz|^2 \le \rhph(\whl) + \lambda_c|\whlc|^2 + \lambda_x\|\whlx\|^2 + \lambda_z|\whlz|^2 \le \rhph(0) + 0 = \ln 2.
\]
Thus,
\[
|\whlz| \le \sqrt{\frac{\ln 2}{\lambda_z}}.
\]
Next, we apply the first-order optimality condition to get 
$$\frac{\partial}{\partial w_z}\Big(\rhph(\whl)  + \lambda_c|\whlc|^2 + \lambda_x\|\whlx\|^2 + \lambda_z|\whlz|^2\Big) = 0 \implies 2\lamz|\whlz| = \Big|\frac{\partial \rhph(\whl)}{\partial \whlz}\Big|.$$ 
Simplifying it and using triangle ineqaulity we get
\[
2\lambda_z |\whlz| = \Big|\frac{\partial \rhph(\whl)}{\partial \whlz}\Big| \le \frac{1}{n}\sum_{i=1}^n |y_i\ph(x_i)\big(1 - \sigma(y_i(\whlc + \whlx^\top x_i + \whlz\ph(x_i)))\big)\Big|.
\]
Since $y_i \in \{-1,1\}$ and $\ph(x_i) \in \{-1,1\}$, we have $|y_i\ph(x_i)| = 1$. Moreover, because the sigmoid function satisfies $0 < \sigma(\cdot) < 1$, we have $0 < 1 - \sigma(\cdot) < 1$. Therefore,
\[
2\lambda_z|\whlz| \le \frac{1}{n}\sum_{i=1}^n 1 = 1 \implies |\whlz| \le \frac{1}{2\lambda_z}.
\]
An identical argument works for $q = \wsl$. Combining both bounds gives
\[
|q_z| \le \min\left\{\sqrt{\frac{\ln 2}{\lambda_z}}, \frac{1}{2\lambda_z}\right\} \le \frac{\ln(2)}{\lamz}.
\]
\end{proof}

The following lemma is an application of McDiarmid's inequality to $\rhph(w)$ and $\rph(w)$.

\begin{lem}\label{lem: McDiamird's inequality}
Let $\ph$ be defined in \eqref{eqn: phi star and phi hat}.
Conditional on $S$, fix $w=(w_c,w_x,w_z)\in\R^{d+2}$. For $0<\delta<1$, with conditional probability at least
$1-\delta$ over $D$,
\begin{equation}
|\rhph(w)-\rph(w)|
\leq(\ln2+\sqrt{\kappa^2+2}\|w\|)
       \sqrt{\frac{2\ln(2/\delta)}{n}}.
\end{equation}
\end{lem}

\begin{proof}
Conditional on $S$, the imputer $\ph$ is fixed and independent of $D$. For any observation $(x_i, y_i) \in D$, let $A_i = (1, x_i, \ph(x_i)) \in \mathbb{R}^{d+2}$. By Assumption (C1) in \eqref{eqn: distribution conditions}, $\|x_i\| \le \kappa$ almost surely, and since $\ph(x_i) \in \{-1, 1\}$, we have
\[
\|A_i\| = \sqrt{1 + \|x_i\|^2 + |\ph(x_i)|^2} \le \sqrt{\kappa^2 + 2}.
\]
By the Cauchy--Schwarz inequality,
\[
|w_c + w_x^\top x_i + w_z\ph(x_i)| = |w^\top A_i| \le \|w\| \|A_i\| \le \sqrt{\kappa^2 + 2}\|w\|.
\]
Since the logistic loss $\ell(x_i, \ph(x_i), y_i; w) = \ln(1 + \exp(-y_i w^\top A_i))$ is $1$-Lipschitz with $\ell(x_i, \ph(x_i), y_i; 0) = \ln 2$, and $y_i \in \{-1, 1\}$, each loss is non-negative and bounded almost surely by
\[
0 \le \ell(x_i, \ph(x_i), y_i; w) \le \ln 2 + |w^\top A_i| \le \ln 2 + \sqrt{\kappa^2 + 2}\|w\| =: M_w.
\]
Replacing any single observation in $D$ changes the empirical risk $\rhph(w) = \frac{1}{n} \sum_{i=1}^n \ell(x_i, \ph(x_i), y_i; w)$ by at most $M_w/n$. Since $\mathbb{E}_D[\rhph(w) \mid S] = \rph(w)$, applying McDiarmid's bounded differences inequality yields, for any $\epsilon > 0$,
\[
\mathbb{P}\big( |\rhph(w) - \rph(w)| \ge \epsilon \mid S \big) \le 2\exp\left(-\frac{2\epsilon^2}{n(M_w/n)^2}\right) = 2\exp\left(-\frac{2n\epsilon^2}{M_w^2}\right).
\]
Setting the right-hand side to $\delta$ and solving for $\epsilon$ gives
\[
\epsilon = M_w \sqrt{\frac{\ln(2/\delta)}{2n}}.
\]
Thus, with conditional probability at least $1 - \delta$ over $D$,
\[
|\rhph(w) - \rph(w)| \le M_w \sqrt{\frac{\ln(2/\delta)}{2n}} \le M_w \sqrt{\frac{2\ln(2/\delta)}{n}} = (\ln 2 + \sqrt{\kappa^2 + 2}\|w\|)\sqrt{\frac{2\ln(2/\delta)}{n}},
\]
which proves the claimed inequality.
\end{proof}


The previous lemma gave a pointwise bound, we now specify a uniform bound.
\begin{lem}\label{lem: uniform ball deviation}
Let $B>0$, $K=\sqrt{\kappa^2+2}$, and define for any $\eta \in (0,1)$:
\[
\Gamma(B,K,\eta):=4KB+(\ln2+KB)\sqrt{\frac{\ln(2/\eta)}2}.
\]
Conditional on $S$, with probability at least $1-\delta$ over $D$,
\[
\sup_{\|w\|\leq B}|\rhph(w)-\rph(w)|
\leq\frac{\Gamma(B,K,\delta)}{\sqrt n}.
\]
\end{lem}
\begin{rem}
In our application, the radius $B$ will be set to the deterministic population envelope $\mathsf{W}$ (defined in Proposition \ref{cor: uniform bound on weights}), ensuring that this uniform concentration bound is applied to a fixed, non-random hypothesis class.
\end{rem}
\begin{proof}
Conditional on $S$, $\ph$ is fixed. Without loss of generality, we can center the loss by subtracting the constant $\ell(x,\ph(x),y;0)=\ln2$ without affecting the difference $\rhph(w)-\rph(w)$. Using standard symmetrization with independent Rademacher variables $\sigma_1,\dots,\sigma_n\in\{-1,1\}$ as shown in \cite[Thm 26.2]{shalev-shwartz_ben-david_2014}, we get
\[
\mathbb{E}_D\left[\sup_{\|w\|\leq B}|\rhph(w)-\rph(w)| \;\middle|\; S\right]
\leq 2\,\mathbb{E}_{D,\sigma}\left[\sup_{\|w\|\leq B}\left|\frac{1}{n}\sum_{i=1}^n\sigma_i\big(\ell(x_i,\ph(x_i),y_i;w)-\ln2\big)\right|\;\middle|\; S\right].
\]
Since the scalar loss $t\mapsto\ln(1+e^{-t})-\ln2$ is $1$-Lipschitz and vanishes at zero, the contraction principle combined with the Cauchy--Schwarz inequality yields
\[
\mathbb{E}_\sigma\left[\sup_{\|w\|\leq B}\left|\frac{1}{n}\sum_{i=1}^n\sigma_i\big(\ell(x_i,\ph(x_i),y_i;w)-\ln2\big)\right|\right]
\leq \frac{2B}{n}\mathbb{E}_\sigma\left\|\sum_{i=1}^n\sigma_i A_i\right\|
\leq \frac{2KB}{\sqrt{n}},
\]
where $A_i=(1,x_i,\ph(x_i))$ satisfies $\|A_i\|\leq\sqrt{\kappa^2+2}=K$ under Assumption (C1), and the final inequality follows from Jensen's inequality with $\mathbb{E}_\sigma\|\sum_{i=1}^n\sigma_i A_i\|^2=\sum_{i=1}^n\|A_i\|^2\leq nK^2$. Hence, the expected uniform deviation is bounded by $\frac{4KB}{\sqrt{n}}$.
To establish concentration around the expectation, define $\Phi(D):=\sup_{\|w\|\leq B}|\rhph(w)-\rph(w)|$. From Lemma \ref{lem: McDiamird's inequality}, each loss on the ball $\|w\|\leq B$ lies in $[0,\ln2+KB]$. Thus, replacing any single observation in $D$ changes $\Phi(D)$ by at most $c_k=\frac{\ln2+KB}{n}$. Applying McDiarmid's bounded differences inequality, the deviation from its expectation satisfies, with conditional probability at least $1-\delta$,
\[
\Phi(D) - \mathbb{E}_D[\Phi(D)\mid S] \leq \frac{\ln2+KB}{\sqrt{n}}\sqrt{\frac{\ln(2/\delta)}{2}}.
\]
Combining the expectation bound with this concentration bound yields
\[
\sup_{\|w\|\leq B}|\rhph(w)-\rph(w)| \leq \frac{4KB}{\sqrt{n}} + \frac{\ln2+KB}{\sqrt{n}}\sqrt{\frac{\ln(2/\delta)}{2}} = \frac{\Gamma(B,K,\delta)}{\sqrt{n}},
\]
which completes the proof.
\end{proof}

The next lemma examines how the optimal weights and risk behave as the regularization parameter varies.

\begin{lem}\label{lem: inverse relation between regulariser and risk}
Let $\lambda=(\lambda_c,\lambda_x,\lambda_z)$ and
$\theta=(\lambda_c,\lambda_x,\theta_z)$ with
$\lambda_z\leq\theta_z$. Then
\begin{equation}\label{eqn: inverse relation between regulariser and risk z weight}
|w^\star_{\theta,z}|\leq|\wslz|.
\end{equation}
If $\lambda_c=\lambda_x=\lambda_0=n^{-1/2}$, then
\begin{equation}\label{eqn: inverse relation between regulariser and risk z}
R(\wsl)\leq R(w^\star_\theta)
+\lambda_0\bigl(|w^\star_{\theta,c}|^2+\|w^\star_{\theta,x}\|^2\bigr).
\end{equation}
The same statements hold for $\rps(\cdot)$ and its corresponding optimal weights.
\end{lem}

\begin{proof}
Write $P(q)=\lambda_c|q_c|^2+\lambda_x\|q_x\|^2$. By the optimality of $\wsl$ for the regularized objective with parameter $\lambda$,
\[
R(\wsl)+P(\wsl)+\lambda_z|\wslz|^2 \leq R(w^\star_\theta)+P(w^\star_\theta)+\lambda_z|w^\star_{\theta,z}|^2.
\]
Similarly, by the optimality of $w^\star_\theta$ for the objective with parameter $\theta$,
\[
R(w^\star_\theta)+P(w^\star_\theta)+\theta_z|w^\star_{\theta,z}|^2 \leq R(\wsl)+P(\wsl)+\theta_z|\wslz|^2.
\]
Adding these two inequalities and canceling common terms yields
\begin{equation}\label{eqn: w theta less than w lambda}
(\theta_z-\lambda_z)\bigl(|w^\star_{\theta,z}|^2-|\wslz|^2\bigr) \leq 0.
\end{equation}
Since $\theta_z \geq \lambda_z > 0$, this implies $|w^\star_{\theta,z}| \leq |\wslz|$, which proves \eqref{eqn: inverse relation between regulariser and risk z weight}.

Using $|w^\star_{\theta,z}|^2 \leq |\wslz|^2$ in the first optimality inequality gives
\[
R(\wsl)+P(\wsl) \leq R(w^\star_\theta)+P(w^\star_\theta) + \lambda_z\big(|w^\star_{\theta,z}|^2 - |\wslz|^2\big) \leq R(w^\star_\theta)+P(w^\star_\theta).
\]
Since $P(\wsl) \geq 0$, we have $R(\wsl) \leq R(\wsl) + P(\wsl)$. Setting $\lambda_c = \lambda_x = \lambda_0 = n^{-1/2}$, we obtain $P(w^\star_\theta) = \lambda_0(|w^\star_{\theta,c}|^2 + \|w^\star_{\theta,x}\|^2)$, which proves \eqref{eqn: inverse relation between regulariser and risk z}. The same argument holds verbatim for $\rps$ by replacing $R$ with $\rps$ and using the corresponding regularized minimizers.
\end{proof}

We next present a corollary which relates the risk of various optimal weights.
\begin{cor}\label{cor: v star and w star infinity}
Let $\lambda=(n^{-1/2},n^{-1/2},\lambda_z)$ with $\lambda_z>0$ and set
$J_\lambda(w)=n^{-1/2}(|w_c|^2+\|w_x\|^2)+\lambda_z|w_z|^2$.
Then
\begin{equation}\label{eqn: population comparator bound}
R(\wsl)+J_\lambda(\wsl)
\leq\min\left\{
R(\ws)+J_\lambda(\ws),\
T(\vs)+\frac{\|\vs\|^2}{\sqrt n}\right\}.
\end{equation}
In particular,
$R(\wsl)\leq T(\vs)+\|\vs\|^2/\sqrt n$.
\end{cor}
\begin{proof}
By the optimality of $\wsl$, we have  
\[
R(\wsl)+J_\lambda(\wsl)\leq R(\ws)+J_\lambda(\ws).
\]
Similarly, evaluating at the feature-only candidate $w=(v_c^\star,v_x^\star,0)\in\mathbb{R}^{d+2}$ yields $R(v_c^\star,v_x^\star,0)=T(\vs)$ and $J_\lambda(v_c^\star,v_x^\star,0)=\frac{\|\vs\|^2}{\sqrt{n}}$, which gives
\[
R(\wsl)+J_\lambda(\wsl)\leq T(\vs)+\frac{\|\vs\|^2}{\sqrt{n}}.
\]
Taking the minimum of these two upper bounds establishes \eqref{eqn: population comparator bound}. Since $J_\lambda(\wsl)\ge 0$, we further get $R(\wsl)\leq T(\vs)+\|\vs\|^2/\sqrt{n}$.
\end{proof}
